\chapter{Reproducibilidad}
\label{anexo:reproducibilidad}

\section{Identidad de los estudios}

La evidencia empírica de este trabajo procede de cuatro estudios encadenados: tres \textit{Model
Studies} en los que el ganador de cada pasada es el \textit{baseline} de la siguiente, y un
\textit{Portfolio Study} que optimiza la cartera del último sin reentrenar nada. Sus identificadores
permiten localizar cualquier cifra citada en el texto.

\begin{description}[leftmargin=3.2cm,style=nextline,font=\normalfont\ttfamily]
  \item[Model Study 1] \studyidA{} · ganador \artefacto{run-803c214fd2ac}
  \item[Model Study 2] \studyidB{} · ganador \artefacto{run-95dcffb1640f}
  \item[Model Study 3] \studyidC{} · ganador \artefacto{run-1bb331e23a18}
  \item[Portfolio Study] \studyportfolio{} · 1.440 carteras evaluadas
\end{description}

El Model Study 3 es el estudio de referencia del documento: de él sale toda la evidencia predictiva.
Sus identificadores de datos son los siguientes.

\begin{description}[leftmargin=3.2cm,style=nextline,font=\normalfont\ttfamily]
  \item[dataset\_hash] \artefacto{b5db6211d03fcd56\ldots{}f182071e5887cd9}
  \item[evaluation\_key] \artefacto{f7152472aeddb96c\ldots{}ac3933d8eb36167d}
  \item[catalog\_hash] \artefacto{930f1f8c106d3928\ldots{}9b9c48857e31910152}
\end{description}

Las cuatro ejecuciones emplearon el mismo catálogo cerrado y la misma ventana de selección --- 117
cohortes mensuales entre 2015 y 2024 ---, y reservaron 2025--2026 sin que interviniera en ninguna
decisión. La métrica de selección fue \texttt{rank\_ic\_only} en las tres pasadas predictivas y
\texttt{information\_ratio} en el Portfolio Study, cuya rejilla se calculó además sobre una serie
recortada en 2024 para que ninguna de las 1.440 combinaciones pudiera observar la era reservada.

\section{Ciclo de vida de una ejecución}

La reproducibilidad no depende sólo de fijar una semilla. Una ejecución puede fallar tras horas de
cálculo, cancelarse o reanudarse, y si esos estados no forman parte del diseño, dos resultados
nominalmente iguales pueden haber recorrido historias distintas. El sistema separa por eso el
\emph{Study}, que contiene la pregunta y el catálogo, de cada \emph{run}, que contiene una
evaluación concreta.

El registro existe antes del cálculo: Study y run se crean antes de lanzar el proceso, y desde ahí
cada transición relevante ---inicio de fase, progreso, finalización, error y cancelación--- queda
escrita. Reanudar consulta primero el estado persistido, de modo que no se repiten ejecuciones
completas ni se duplican decisiones ya tomadas.

Aislar cada evaluación en su propio proceso tiene además una función científica: evita que un estado
en memoria ---una semilla, una caché, un dataframe mutado--- sobreviva de una configuración a la
siguiente. El resultado depende así de entradas persistidas y no del orden accidental en que se
ejecutaron las cosas.

\section{Cadena de custodia de la evidencia}

Cada nivel responde a una pregunta diferente y tiene un consumidor definido:

\begin{table}[H]
  \centering
  \small
  \begin{tabularx}{\textwidth}{p{3.1cm}X X}
    \toprule
    Nivel & Contenido & Función en la auditoría \\
    \midrule
    Catálogo & Rejillas, restricciones y fases & Demuestra qué podía elegirse antes de medir. \\
    \texttt{winner.json} & Configuración adoptada y procedencia & Une cada valor final con el run que lo produjo. \\
    \texttt{decisions.json} & Incumbente, candidatos y puerta pareada & Reconstruye por qué cambió o no cada decisión. \\
    Evidencia predictiva & Scores, Rank-IC, atribución y robustez & Sostiene las afirmaciones sobre ordenación. \\
    Evidencia de cartera & Órdenes, posiciones, costes y rejilla & Sostiene las afirmaciones económicas. \\
    Manifiesto LaTeX & Fuente de cada tabla, macro y figura & Impide que el manuscrito mezcle estudios. \\
    \bottomrule
  \end{tabularx}
  \caption{Niveles de persistencia y propósito de cada uno.}
  \label{tab:cadena-custodia}
\end{table}

La identidad del panel señala qué datos se usaron, la del catálogo qué espacio de decisiones estaba
disponible, y la clave de evaluación une ambas con la configuración concreta. Una caché sólo es
reutilizable cuando las tres coinciden y el artefacto está completo, lo que evita dos formas de
contaminación difíciles de detectar: atribuir a un catálogo nuevo un resultado antiguo, y dar por
terminado un directorio que quedó a medias tras una interrupción.

\section{Regeneración de figuras y tablas}

Las figuras y tablas del manuscrito no se dibujan a mano: se generan desde los artefactos
persistidos mediante un único comando, ejecutado desde la raíz del repositorio.

\begin{quote}
\ttfamily
python latex/scripts/export\_study\_assets.py \\
\hspace*{2em}--study-id study-20260817-094411-568bd37e \\
\hspace*{2em}--chain-study-id study-20260816-182345-3cc1a5fb \\
\hspace*{2em}--chain-study-id study-20260817-021135-b5926b62 \\
\hspace*{2em}--chain-study-id study-20260817-094411-568bd37e \\
\hspace*{2em}--portfolio-study-id study-20260817-212856-f86ca822
\end{quote}

Los tres identificadores de cadena van en orden de ejecución y el último coincide con el estudio de
referencia. La separación entre el estudio predictivo y el de cartera hace explícita la procedencia
de cada activo: del primero sale todo lo predictivo ---Rank-IC, agentes, meta-agente, decisiones,
robustez y atribución--- y del segundo todo lo económico ---rejilla de carteras, patrimonio,
órdenes, métricas anuales y perfiles---.

El script escribe las figuras y los cuerpos de tabla en sus carpetas, genera el fichero de macros
numéricas que consumen la memoria y la defensa, y actualiza el manifiesto de activos con la lista de
artefactos consumidos, separando las fuentes predictivas de las económicas. Ese manifiesto es la
tabla de correspondencia entre cada salida del manuscrito y su origen: ninguna cifra generada por
script entra en el documento sin quedar registrada allí, y por eso no se imprime una ruta junto a
cada cifra. La única excepción es la Tabla~\ref{tab:defectos} del Anexo~\ref{anexo:desarrollo}, que
se transcribe de la bitácora del proyecto porque documenta el desarrollo y no una medición.

Una salvedad sobre qué se puede regenerar. El repositorio versiona los artefactos en formato JSON,
que sostienen la mayoría de las cifras citadas y permiten verificarlas directamente; los ficheros
Parquet quedan excluidos por tamaño, de modo que las tablas y figuras construidas a partir de ellos
se conservan ya generadas pero no pueden recalcularse desde un clon del repositorio sin volver a
ejecutar el estudio.

Otros dos comandos completan la verificación: uno comprueba que el proyecto sea portable ---rutas
relativas, recursos existentes, referencias cruzadas con destino, ausencia de \textit{mojibake} y de
activos huérfanos--- y otro recalcula en memoria las macros decisivas, sin escribir nada, y falla si
difieren de los estudios adoptados.

